We have developed a framework for compilation and optimization of application scale quantum circuits that can negotiate the trade off between reducing resource requirements and providing rigorous guarantees on output error. We demonstrated this framework across a diverse set of benchmark quantum algorithm circuits, using specific tools (\eg BQSKit, GridSynth) for components of our framework, and showed consistent reductions in resource-intensive gates -- CNOTs for NISQ architectures and $T$ gates for FT architectures -- while simultaneously controlling output errors.  
For benchmarks where the weighted ensemble error is quadratically reduced (\ie scales as $O(\epsilon^2)$, when circuits composing the ensemble are synthesized to approximation error $\epsilon$) for a significant fraction of circuit blocks, we observe resource reductions exceeding $40\%$ in some cases, without exceeding the prescribed error tolerance.  
Output state error and output observable errors scale as $O(\epsilon^2)$, in agreement with theoretical predictions, confirming that the quadratic error suppression extends from the block level to the full compiled circuit.  
These results demonstrate that ensemble-based approximate compilation can simultaneously deliver meaningful resource savings and provable error guarantees, making it a viable strategy for both near-term and fault-tolerant quantum computing regimes.

Our demonstrations used specific, extant tools to implement components of our workflow such as circuit partitioning and circuit synthesis. Our numerics validate the performance of these implementations. However, we emphasize that the overall framework, summarized in Fig. \ref{fig:resut_workflow}, and its theoretical foundations are the central contribution of this work. This framework allows for future development and maturation of components, especially the circuit synthesis and diversification components, that might allow for even better negotiation of the trade-off between resource reduction and control of output errors. 